\section{Related literature}
\label{sec:related}

\subsection{Continuous approximations and spatial service design}

\citet{Beardwood1959} establish an asymptotic relation between Euclidean tour length, the number of stops, and their spatial distribution. Partitioning algorithms \citep{Karp1977} and analytical fleet-routing models with stop limits \citep{Daganzo1984} connect this geometric scaling to spatial service organization. Our capacity tables use a BHH-inspired approximation with handling and stem travel specified separately. The equal division of wave routing work among vehicles remains a finite-tour modeling approximation.

\citet{Banerjee2022} use continuous approximations to maximize a single vehicle's service-zone area, characterize the associated dispatch policies, and calculate maximum-area functions. Those functions provide fleet-size guidelines and inputs to approximate weighted Voronoi partitioning. Their framework separates depot--zone travel, routing within the zone, setup, and per-order service time. Their term \emph{linehaul} denotes the depot--zone connection; here, line-haul denotes intercity hub-to-hub transport, and local connections are called stem travel. The decision levels also differ: they design single-vehicle zones and fleet sizes, whereas we schedule a given finite fleet across two cities. Our within-wave balanced partitioning assumption does not determine persistent vehicle territories.

\subsection{Consolidation and dispatch decisions}

Transportation and inventory costs jointly shape consolidation decisions \citep{Burns1985}; terminals further change the economics of many-to-many freight flows \citep{Daganzo1987}. \citet{CetinkayaBookbinder2003} analyze quantity-based and time-based shipment consolidation through renewal costs. \citet{Wei2023} compare hybrid, time-based, and quantity-based policies with inbound replenishment and outbound dispatch. These studies motivate explicit separation of batch-service cost and accumulation waiting. Coordinated cycles also arise in integer-ratio inventory policies \citep{Roundy1985}, whose performance guarantees depend on their inventory setting.

Departure timing has long been a decision variable: \citet{Newell1971} studies dispatch under smooth time-varying arrivals, and \citet{Klapp2018} formulate a same-day delivery problem in which a vehicle can dispatch ready orders or wait at available decision epochs. Their dispatch waves are decision opportunities; our local waves identify service intervals whose pooled stops determine capacity. Stochastic routing studies additionally retain random arrival locations and finite vehicle resources \citep{Bertsimas1991,Bertsimas1993}. Those features distinguish full routing policies from a scalar backlog-reset model.

Our stationary benchmark uses a prescribed common-wave class and a routing-based service cost optimized over local vehicle loads. Its contribution is a verified connection between that cost and admissible repeated dispatch decisions. Impulse-control QVIs and verification are established methods \citep{LionsPerthame1986,Belak2017}; the appendices give a self-contained verification for the deterministic separable benchmark. The Poisson comparison concerns the same positive integer full-reset thresholds, without asserting optimality over arbitrary stochastic policies.

\subsection{Multistage freight scheduling and vehicle circulation}

Time-expanded service networks connect departure decisions to resource availability. \citet{Scherr2020} study two-tier logistics with mixed autonomous fleets and develop dynamic discretization discovery for the resulting service network. \citet{He2022} develop an exact algorithm for service network design with hub capacity constraints. Their methods address time discretization and network coupling; our formulation instead uses fixed time nodes and precomputed wave capacities.

\citet{Lee2024} optimize autonomous transfer hub networks combining automated middle-mile transport and human-driven endpoint service. We retain that fleet separation while representing local service through wave capacities and explicitly accounting for human-driven direct and empty movements between cities. Cumulative city inventories make relocation affect subsequent local availability, and cargo precedence connects collection, line-haul, and distribution. Rolling-horizon execution applies this model to revealed orders and preserves earlier commitments. Candidate filtering and independent replay support its implementation, with their scope stated alongside the procedure.
